\documentclass[../../main.tex]{subfiles}

\begin{document}

\section{欧氏空间 $\mathbf{R}^{n}$ 上的 Lebesgue 测度}\label{section:ch20-lebesgue}

对自然数 $n$, 我们记 $n$ 元有序实数组 $x=(x_{1},\cdots,x_{n})$ 全体为 $\mathbf{R}^{n}$. 则 $\mathbf{R}^{n}$ 是线性空间且存在如下定义的双线性形式 $\langle\cdot,\cdot\rangle:\mathbf{R}^{n}\times\mathbf{R}^{n}\to\mathbf{R}$:
\begin{align*}
\text{对所有 }x,y\in\mathbf{R}^{n},\quad\langle x,y\rangle=\sum_{k=1}^{n}x_{k}\cdot y_{k}.
\end{align*}
该双线性形式称为内积或数量积. 它诱导了 $\mathbf{R}^{n}$ 上的范数:
\begin{align*}
\text{对所有 }x\in\mathbf{R}^{n},\quad\|x\|=\sqrt{\langle x,x\rangle}=\sqrt{\sum_{k=1}^{n}x_{k}^{2}}.
\end{align*}
从而诱导了 $\mathbf{R}^{n}$ 上的拓扑. 该范数称为欧氏范数. 具有该内积与它所诱导的范数和拓扑的线性空间 $\mathbf{R}^{n}$ 称为 $n$ 维欧氏空间.

谈到 $\mathbf{R}$ 中的有界区间, 我们意味着形如 $[a,b]$, $[a,b)$, $(a,b]$ 或 $(a,b)$ 的集合, 这里实数 $a\leqslant b$. 我们考虑空集和由单点组成的集合为有界区间. 对于端点为 $a$ 与 $b$ 的有界区间 $I$, 我们定义它的长度 $\ell(I)$ 为 $b-a$.

\begin{definition}[$\mathbf{R}^{n}$ 中的有界区间]\label{definition:ch20-有界区间}
谈到 $\mathbf{R}^{n}$ 中的有界区间, 我们指由 $n$ 个实数的有界区间的笛卡儿积构成的集合 $I$,
\begin{align*}
I=I_{1}\times I_{2}\times\cdots\times I_{n}.
\end{align*}
定义 $I$ 的体积 $\operatorname{vol}(I)$ 为
\begin{align*}
\operatorname{vol}(I)=\ell(I_{1})\cdot\ell(I_{2})\cdots\ell(I_{n}).
\end{align*}
\end{definition}

\begin{definition}[整点计数]\label{definition:ch20-整点计数}
我们称 $\mathbf{R}^{n}$ 中的点为整点, 若它的每个坐标都是整数. 对 $\mathbf{R}^{n}$ 中的有界区间 $I$, 我们定义它的整点计数 $\mu^{\operatorname{integral}}(I)$ 为 $I$ 中的整点个数.
\end{definition}

\begin{lemma}\label{lemma:ch20-整点计数极限}
对每个 $\varepsilon>0$, 定义扩张 $T_{\varepsilon}:\mathbf{R}^{n}\to\mathbf{R}^{n}$ 为 $T_{\varepsilon}(x)=\varepsilon\cdot x$. 则对 $\mathbf{R}^{n}$ 中的每个有界区间 $I$,
\begin{align}\label{align:ch20-整点极限}
\lim_{\varepsilon\to\infty}\frac{\mu^{\operatorname{integral}}(T_{\varepsilon}(I))}{\varepsilon^{n}}=\operatorname{vol}(I).
\end{align}
\end{lemma}

\begin{proof}
对于 $\mathbf{R}$ 中具有端点 $a$ 和 $b$ 的有界区间 $I$, 关于 $I$ 的整点计数, 我们有以下估计(见习题 18):
\begin{align}\label{align:ch20-一维整点估计}
(b-a)-1\leqslant\mu^{\operatorname{integral}}(I)\leqslant(b-a)+1.
\end{align}
因此对于区间 $I=I_{1}\times I_{2}\times\cdots\times I_{n}$, 由于
\begin{align*}
\mu^{\operatorname{integral}}(I)=\mu^{\operatorname{integral}}(I_{1})\cdots\mu^{\operatorname{integral}}(I_{n}),
\end{align*}
若每个区间 $I_{k}$ 具有端点 $a_{k}$ 与 $b_{k}$, 我们有估计
\begin{align}\label{align:ch20-n维整点估计}
[(b_{1}-a_{1})-1]\cdots[(b_{n}-a_{n})-1]\leqslant\mu^{\operatorname{integral}}(I)\leqslant[(b_{1}-a_{1})+1]\cdots[(b_{n}-a_{n})+1].
\end{align}
对 $\varepsilon>0$, 我们用扩张区间 $T_{\varepsilon}(I)$ 代替区间 $I$ 得到估计
\begin{align}\label{align:ch20-扩张整点估计}
[\varepsilon\cdot(b_{1}-a_{1})-1]\cdots[\varepsilon\cdot(b_{n}-a_{n})-1]&\leqslant\mu^{\operatorname{integral}}(T_{\varepsilon}(I))\\
&\leqslant[\varepsilon\cdot(b_{1}-a_{1})+1]\cdots[\varepsilon\cdot(b_{n}-a_{n})+1].\nonumber
\end{align}
不等式除以 $\varepsilon^{n}$ 且取 $\varepsilon\to\infty$ 的极限我们得到\eqref{align:ch20-整点极限}.
\end{proof}

我们将以下命题的证明留作归纳法的练习, 证明中用到两个半环的笛卡儿积是半环这一性质(见习题 25).

\begin{proposition}\label{proposition:ch20-区间半环}
$\mathbf{R}^{n}$ 中的有界区间族 $\mathcal{I}$ 是半环.
\end{proposition}

\begin{proposition}\label{proposition:ch20-体积预测度}
体积集函数 $\operatorname{vol}:\mathcal{I}\to[0,\infty]$ 是 $\mathbf{R}^{n}$ 中的有界区间构成的半环 $\mathcal{I}$ 上的预测度.
\end{proposition}

\begin{proof}
我们首先证明体积在有界区间的有限不交并上是有限可加的. 令 $I$ 为 $\mathbf{R}^{n}$ 中的有界区间, 它是有界区间的有限不交有限族 $\{I^{k}\}_{k=1}^{m}$ 的并. 则对每个 $\varepsilon>0$, 有界区间 $T_{\varepsilon}(I)$ 是有限个不交有界区间 $\{T_{\varepsilon}(I^{k})\}_{k=1}^{m}$ 的并. 显然 $\mu^{\operatorname{integral}}$ 是有限可加的. 因此,
\begin{align*}
\text{对所有 }\varepsilon>0,\quad\mu^{\operatorname{integral}}(T_{\varepsilon}(I))=\sum_{k=1}^{m}\mu^{\operatorname{integral}}(T_{\varepsilon}(I^{k})).
\end{align*}
每边除以 $\varepsilon^{n}$ 取 $\varepsilon\to\infty$ 的极限, 由\eqref{align:ch20-整点极限}得到
\begin{align*}
\operatorname{vol}(I)=\sum_{k=1}^{m}\operatorname{vol}(I^{k}).
\end{align*}
因此体积是有限可加的.

剩下来要证明体积的可数单调性. 令 $I$ 为 $\mathbf{R}^{n}$ 中的有界区间, 它被可数个有界区间 $\{I^{k}\}_{k=1}^{\infty}$ 覆盖. 我们首先考虑 $I$ 是闭区间且每个 $I^{k}$ 是开的情形. 根据 Heine-Borel 定理, 我们可以选取自然数 $m$ 使得 $I$ 可被有限子族 $\{I^{k}\}_{k=1}^{m}$ 覆盖. 显然整点计数 $\mu^{\operatorname{integral}}$ 是单调与有限可加的, 因为区间族是半环、有限单调的. 因此,
\begin{align*}
\mu^{\operatorname{integral}}(I)\leqslant\sum_{k=1}^{m}\mu^{\operatorname{integral}}(I^{k}).
\end{align*}
扩张这些区间. 因此,
\begin{align*}
\text{对所有 }\varepsilon>0,\quad\mu^{\operatorname{integral}}(T_{\varepsilon}(I))\leqslant\sum_{k=1}^{m}\mu^{\operatorname{integral}}(T_{\varepsilon}(I^{k})).
\end{align*}
每边除以 $\varepsilon^{n}$ 取 $\varepsilon\to\infty$ 的极限, 由\eqref{align:ch20-整点极限}得到
\begin{align*}
\operatorname{vol}(I)\leqslant\sum_{k=1}^{m}\operatorname{vol}(I^{k})\leqslant\sum_{k=1}^{\infty}\operatorname{vol}(I^{k}).
\end{align*}
剩下来要考虑覆盖区间 $I$ 的一般有界区间族 $\{I^{k}\}_{k=1}^{\infty}$ 的情形. 令 $\varepsilon>0$. 选取闭区间 $\widehat{I}$ 包含于 $I$ 与开区间族 $\{\widehat{I}^{k}\}_{k=1}^{\infty}$ 使得每个 $I^{m}\subseteq\widehat{I}^{m}$, 此外,
\begin{align*}
\text{对所有 }m,\operatorname{vol}(I)-\operatorname{vol}(\widehat{I})<\varepsilon\text{ 且 }\operatorname{vol}(\widehat{I}^{m})-\operatorname{vol}(I^{m})<\varepsilon/2^{m}.
\end{align*}
根据刚考虑过的情形,
\begin{align*}
\operatorname{vol}(\widehat{I})\leqslant\sum_{k=1}^{\infty}\operatorname{vol}(\widehat{I}^{k}).
\end{align*}
因此,
\begin{align*}
\operatorname{vol}(I)\leqslant\sum_{k=1}^{\infty}\operatorname{vol}(I^{k})+2\varepsilon.
\end{align*}
由于这对所有 $\varepsilon>0$ 成立, 它对 $\varepsilon=0$ 也成立. 因此体积集函数是预测度.
\end{proof}

\begin{definition}[Lebesgue 外测度与 Lebesgue 测度]\label{definition:ch20-lebesgue测度}
由 $\mathbf{R}^{n}$ 中的有界区间的半环上的体积预测度诱导的外测度 $\mu_{n}^{*}$ 称为 $\mathbf{R}^{n}$ 上的 Lebesgue 外测度. $\mu_{n}^{*}$ 可测集族记为 $\mathcal{L}^{n}$, 该族中的集合称为 Lebesgue 可测集. $\mu_{n}^{*}$ 在 $\mathcal{L}^{n}$ 上的限制称为 $\mathbf{R}^{n}$ 上的 Lebesgue 测度或 $n$ 维 Lebesgue 测度, 记为 $\mu_{n}$.
\end{definition}

\begin{theorem}\label{theorem:ch20-lebesgue基本性质}
$\mathbf{R}^{n}$ 的 Lebesgue 可测子集的 $\sigma$ 代数 $\mathcal{L}^{n}$ 包含 $\mathbf{R}^{n}$ 中的有界区间, 以及 $\mathbf{R}^{n}$ 的 Borel 子集. 此外, 测度空间 $(\mathbf{R}^{n},\mathcal{L}^{n},\mu_{n})$ 是 $\sigma$ 有限与完备的且对 $\mathbf{R}^{n}$ 中的有界区间 $I$,
\begin{align*}
\mu_{n}(I)=\operatorname{vol}(I).
\end{align*}
\end{theorem}

\begin{proof}
根据前一个命题, 体积是 $\mathbf{R}^{n}$ 中的有界区间的半环上的预测度. 它显然是 $\sigma$ 有限的. 因此 Carath\'eodory-Hahn 定理告诉我们, Lebesgue 测度是体积的延拓, 而测度空间 $(\mathbf{R}^{n},\mathcal{L}^{n},\mu_{n})$ 既是 $\sigma$ 有限又是完备的. 剩下来要证明每个 Borel 集是 Lebesgue 可测的. 由于 Borel 集族是包含开集的最小 $\sigma$ 代数, 仅须证明 $\mathbf{R}^{n}$ 的每个开子集 $\mathcal{O}$ 是 Lebesgue 可测的. 令 $\mathcal{O}$ 为 $\mathbf{R}^{n}$ 中的开集. $\mathcal{O}$ 中具有有理坐标的点族是 $\mathcal{O}$ 的可数稠密子集. 令 $\{z_{k}\}_{k=1}^{\infty}$ 为该族的一个列举. 对每个 $k$, 考虑中心在 $z_{k}$、长度为 $1/m$ 的开方体 $I_{k,m}$. 我们把证明
\begin{align}\label{align:ch20-开集方体分解}
\mathcal{O}=\bigcup_{I_{k,m}\subseteq\mathcal{O}}I_{k,m}
\end{align}
留作练习. 因此每个 $I_{k,m}$ 是可测的, 作为这些集的可数并 $\mathcal{O}$ 也是可测的.
\end{proof}

\begin{corollary}\label{corollary:ch20-连续函数可测}
令 $E$ 为 $\mathbf{R}^{n}$ 的 Lebesgue 可测子集, 而 $f:E\to\mathbf{R}$ 是连续的. 则 $f$ 关于 $n$ 维 Lebesgue 测度是可测的.
\end{corollary}

\begin{proof}
令 $\mathcal{O}$ 为实数的开集. 则根据 $f$ 在 $E$ 上的连续性, $f^{-1}(\mathcal{O})=E\cap\mathcal{U}$, 其中 $\mathcal{U}$ 在 $\mathbf{R}^{n}$ 中是开的. 根据前一个定理, $\mathcal{U}$ 是可测的, 因此 $f^{-1}(\mathcal{O})$ 是可测的.
\end{proof}

\textbf{Lebesgue 测度的正则性.} 以下定理和它的推论把 $\mathbf{R}^{n}$ 上的 Lebesgue 测度与 $\mathbf{R}^{n}$ 上由欧氏范数诱导的拓扑紧密地联系起来.

\begin{theorem}\label{theorem:ch20-lebesgue正则}
令 $E$ 为 $\mathbf{R}^{n}$ 的 Lebesgue 可测子集. 则
\begin{align}\label{align:ch20-外正则}
\mu_{n}(E)=\inf\{\mu_{n}(\mathcal{O})\mid E\subseteq\mathcal{O},\mathcal{O}\text{ 是开的}\}
\end{align}
且
\begin{align}\label{align:ch20-内正则}
\mu_{n}(E)=\sup\{\mu_{n}(K)\mid K\subseteq E,K\text{ 是紧的}\}.
\end{align}
\end{theorem}

\begin{proof}
我们考虑 $E$ 是有界的, 因此有有限 Lebesgue 测度的情形, 而将推广到无界的 $E$ 作为练习. 我们首先证明\eqref{align:ch20-外正则}. 令 $\varepsilon>0$. 由于 $\mu_{n}(E)=\mu_{n}^{*}(E)<\infty$, 根据 Lebesgue 外测度的定义, 我们可以选取 $\mathbf{R}^{n}$ 中的有界区间的可数族 $\{I^{m}\}_{m=1}^{\infty}$, 它覆盖 $E$ 且
\begin{align*}
\sum_{m=1}^{\infty}\mu_{n}(I^{m})<\mu_{n}(E)+\varepsilon/2.
\end{align*}
对每个 $m$, 选取包含 $I^{m}$ 且测度小于 $\mu_{n}(I^{m})+\varepsilon/[2^{m+1}]$ 的开区间. 该开区间族的并是开集, 我们把它记为 $\mathcal{O}$. 则 $E\subseteq\mathcal{O}$, 且根据测度的可数单调性, $\mu_{n}(\mathcal{O})<\mu_{n}(E)+\varepsilon$. 因此\eqref{align:ch20-外正则}得证.

我们现在证明\eqref{align:ch20-内正则}. 由于 $E$ 是有界的, 我们可以选取包含 $E$ 的闭有界集 $K'$. 由于 $K'\setminus E$ 是有界的, 我们从证明的第一部分推出存在开集 $\mathcal{O}$ 使得 $K'\setminus E\subseteq\mathcal{O}$, 根据 $\mu_{n}$ 的分割性质,
\begin{align}\label{align:ch20-开集估计}
\mu_{n}(\mathcal{O}\setminus[K'\setminus E])<\varepsilon.
\end{align}
定义
\begin{align*}
K=K'\setminus\mathcal{O}.
\end{align*}
则 $K$ 是闭有界的, 因此是紧的. 从包含关系 $K'\setminus E\subseteq\mathcal{O}$ 与 $E\subseteq K'$ 我们推出,
\begin{align*}
K=K'\setminus\mathcal{O}\subseteq K'\setminus[K'\setminus E]=K'\cap E\subseteq E.
\end{align*}
因此 $K\subseteq E$. 另一方面, 从包含关系 $E\subseteq K'$ 我们推出,
\begin{align*}
E\setminus K=E\setminus[K'\setminus\mathcal{O}]=E\cap\mathcal{O}
\end{align*}
与
\begin{align*}
E\cap\mathcal{O}\subseteq\mathcal{O}\setminus[K'\setminus E].
\end{align*}
因此根据测度的分割性与单调性以及\eqref{align:ch20-开集估计},
\begin{align*}
\mu_{n}(E)-\mu_{n}(K)=\mu_{n}(E\setminus K)\leqslant\mu_{n}(\mathcal{O}\setminus[K'\setminus E])<\varepsilon.
\end{align*}
于是\eqref{align:ch20-内正则}得证. 定理证毕.
\end{proof}

$\mathbf{R}^{n}$ 的每个 Borel 子集是 Lebesgue 可测的, 因此任何 $G_{\delta}$ 或 $F_{\sigma}$ 集也如此. 此外, 每个 Lebesgue 外测度为零的集合是 Lebesgue 可测的. 因此前一个定理以及测度的连续性与分割性质, 给出了 Lebesgue 可测集的相对简单的刻画. 它应该与第 17 章的命题 10 相比较.

\begin{corollary}\label{corollary:ch20-可测集刻画}
对 $\mathbf{R}^{n}$ 的子集 $E$, 以下假设是等价的:
\begin{enumerate}[(i)]
\item $E$ 关于 $n$ 维 Lebesgue 测度是可测的.
\item 存在 $\mathbf{R}^{n}$ 的 $G_{\delta}$ 子集 $G$ 使得
\begin{align*}
E\subseteq G\text{ 且 }\mu_{n}^{*}(G\setminus E)=0.
\end{align*}
\item 存在 $\mathbf{R}^{n}$ 的 $F_{\sigma}$ 子集 $F$ 使得
\begin{align*}
F\subseteq E\text{ 且 }\mu_{n}^{*}(E\setminus F)=0.
\end{align*}
\end{enumerate}
\end{corollary}

从上面的 Lebesgue 可测集的刻画推出 Lebesgue 测度是平移不变的: 对 $E\subseteq\mathbf{R}^{n}$ 与 $z\in\mathbf{R}^{n}$, 定义 $E$ 的 $z$ 平移为
\begin{align*}
E+z=\{x+z\mid x\in E\}.
\end{align*}
若 $E$ 是 $\mu_{n}$ 可测的, $E+z$ 也如此, 且
\begin{align*}
\mu_{n}(E)=\mu_{n}(E+z).
\end{align*}

\textbf{作为乘积测度的 Lebesgue 测度.} 对满足 $n=m+k$ 的自然数 $n$, $m$ 与 $k$, 考虑集合 $\mathbf{R}^{n}$、$\mathbf{R}^{m}$、$\mathbf{R}^{k}$ 与 $\mathbf{R}^{m}\times\mathbf{R}^{k}$ 以及定义为
\begin{align}\label{align:ch20-坐标分解}
\text{对所有 }x=(x_{1},\cdots,x_{n})\in\mathbf{R}^{n},\quad\varphi(x_{1},\cdots,x_{n})=((x_{1},\cdots,x_{m}),(x_{m+1},\cdots,x_{m+k}))
\end{align}
的映射
\begin{align*}
\varphi:\mathbf{R}^{n}\to\mathbf{R}^{m}\times\mathbf{R}^{k}.
\end{align*}
该映射是一对一与映上的. 每个集合 $\mathbf{R}^{n}$、$\mathbf{R}^{m}$ 与 $\mathbf{R}^{k}$ 有线性结构、拓扑结构、测度结构, 乘积空间 $\mathbf{R}^{m}\times\mathbf{R}^{k}$ 从成员空间继承了线性结构、拓扑结构与测度结构. 映射 $\varphi$ 关于线性结构与拓扑结构是同构的. 以下命题告诉我们从测度的观点, 映射 $\varphi$ 也是同构.

\begin{proposition}\label{proposition:ch20-乘积同构}
对于由\eqref{align:ch20-坐标分解}定义的映射 $\varphi:\mathbf{R}^{n}\to\mathbf{R}^{m}\times\mathbf{R}^{k}$, $\mathbf{R}^{n}$ 的子集 $E$ 关于 $n$ 维 Lebesgue 测度是可测的当且仅当它的象 $\varphi(E)$ 关于 $\mathbf{R}^{m}\times\mathbf{R}^{k}$ 上的乘积测度 $\mu_{m}\times\mu_{k}$ 是可测的, 且
\begin{align*}
\mu_{n}(E)=(\mu_{m}\times\mu_{k})(\varphi(E)).
\end{align*}
\end{proposition}

\begin{proof}
定义 $\mathcal{I}_{n}$ 为 $\mathbf{R}^{n}$ 中的区间族而 $\operatorname{vol}_{n}$ 是定义在 $\mathcal{I}_{n}$ 上的体积集函数. 由于 $\operatorname{vol}_{n}$ 是 $\sigma$ 有限预测度, 从 Carath\'eodory-Hahn 定理的唯一性部分得出, Lebesgue 测度 $\mu_{n}$ 是 $\mathcal{L}^{n}$ 上的延拓 $\operatorname{vol}_{n}:\mathcal{I}_{n}\to[0,\infty]$ 的唯一测度.

显然,
\begin{align}\label{align:ch20-区间体积保持}
\text{对所有 }I\in\mathcal{I}_{n},\quad\mu_{n}(I)=(\mu_{m}\times\mu_{k})(\varphi(I)).
\end{align}
我们把证明这蕴涵外测度被 $\varphi$ 保持, 因此 $E$ 属于 $\mathcal{L}^{n}$ 当且仅当 $\varphi(E)$ 是 $(\mu_{m}\times\mu_{k})$ 可测的留给读者作为练习. 由于 $\varphi$ 是一对一与映上的, 若我们对所有 $E\in\mathcal{L}^{n}$ 定义
\begin{align*}
\mu'(E)=(\mu_{m}\times\mu_{k})(\varphi(E)),
\end{align*}
则 $\mu'$ 是 $\mathcal{L}^{n}$ 上的延拓 $\operatorname{vol}_{n}:\mathcal{I}_{n}\to[0,\infty]$ 的测度. 因此根据 $\mu_{n}$ 的唯一性断言,
\begin{align*}
\text{对所有 }E\in\mathcal{L}^{n},\quad\mu_{n}(E)=\mu'(E)=(\mu_{m}\times\mu_{k})(\varphi(E)).
\end{align*}
这完成了证明.
\end{proof}

从该命题、Lebesgue 测度的 $\sigma$ 有限性与完备性以及 Fubini 定理与 Tonelli 定理, 我们有关于 $\mathbf{R}^{n}$ 上的 Lebesgue 测度的积分的定理.

\begin{theorem}\label{theorem:ch20-富比尼Lebesgue}
对使得 $n=m+k$ 的自然数 $n$, $m$ 与 $k$, 考虑由\eqref{align:ch20-坐标分解}定义的映射 $\varphi:\mathbf{R}^{n}\to\mathbf{R}^{m}\times\mathbf{R}^{k}$. 函数 $f:\mathbf{R}^{m}\times\mathbf{R}^{k}\to\mathbf{R}$ 关于乘积测度 $\mu_{m}\times\mu_{k}$ 可测当且仅当复合 $f\circ\varphi:\mathbf{R}^{n}\to\mathbf{R}$ 关于 Lebesgue 测度 $\mu_{n}$ 可测. 若 $f$ 在 $\mathbf{R}^{n}$ 上关于 Lebesgue 测度 $\mu_{n}$ 可积, 则
\begin{align}\label{align:ch20-累次积分}
\int_{\mathbf{R}^{n}}f\,d\mu_{n}=\int_{\mathbf{R}^{k}}\left[\int_{\mathbf{R}^{m}}f(x,y)\,d\mu_{m}(x)\right]d\mu_{k}(y).
\end{align}
此外, 若 $f$ 非负且关于 Lebesgue 测度 $\mu_{n}$ 可测, 则上面的等式也成立.
\end{theorem}

\textbf{Lebesgue 积分与线性换元.} 我们记 $L(\mathbf{R}^{n})$ 为线性算子 $T:\mathbf{R}^{n}\to\mathbf{R}^{n}$ 的线性空间, 记 $GL(n,\mathbf{R})$ 为 $L(\mathbf{R}^{n})$ 的那些由可逆线性算子 $T:\mathbf{R}^{n}\to\mathbf{R}^{n}$(即一对一与映上的线性算子)组成的子集. 可逆算子的逆是线性的. 在算子的复合下, $GL(n,\mathbf{R})$ 是一个群, 它称为 $\mathbf{R}^{n}$ 的一般线性群. 对 $1\leqslant k\leqslant n$, 我们把 $\mathbf{R}^{n}$ 中的 $k$ 重坐标是 $1$ 而其他坐标是 $0$ 的点记为 $e_{k}$. 则 $\{e_{1},\cdots,e_{n}\}$ 是 $L(\mathbf{R}^{n})$ 的典范基. 观察到一旦 $T(e_{k})$ 给定, $1\leqslant k\leqslant n$, 则线性算子 $T:\mathbf{R}^{n}\to\mathbf{R}^{n}$ 被唯一确定, 这是由于若 $x=(x_{1},\cdots,x_{n})$, 则
\begin{align*}
\text{对所有 }x\in\mathbf{R}^{n},\quad T(x)=T(x_{1}e_{1}+\cdots+x_{n}e_{n})=x_{1}T(e_{1})+\cdots+x_{n}T(e_{n}).
\end{align*}
我们仅须线性算子的一个分析性质: 它是 Lipschitz 的.

\begin{proposition}\label{proposition:ch20-线性算子lipschitz}
线性算子 $T:\mathbf{R}^{n}\to\mathbf{R}^{n}$ 是 Lipschitz 的.
\end{proposition}

\begin{proof}
令 $x$ 属于 $\mathbf{R}^{n}$. 正如我们已观察到的, 根据 $T$ 的线性,
\begin{align*}
\text{对所有 }x\in\mathbf{R}^{n},\quad T(x)=x_{1}T(e_{1})+\cdots+x_{n}T(e_{n}).
\end{align*}
因此, 根据范数的次可加性与正齐次性,
\begin{align*}
\|T(x)\|=\|x_{1}T(e_{1})+\cdots+x_{n}T(e_{n})\|\leqslant\sum_{k=1}^{n}|x_{k}|\cdot\|T(e_{k})\|.
\end{align*}
因此, 若我们定义 $c=\sqrt{\sum_{k=1}^{n}\|T(e_{k})\|^{2}}$, 根据 Cauchy-Schwarz 不等式,
\begin{align*}
\|T(x)\|\leqslant c\cdot\|x\|.
\end{align*}
对任何 $u,v\in\mathbf{R}^{n}$, 设 $x=u-v$. 则根据 $T$ 的线性, $T(x)=T(u-v)=T(u)-T(v)$, 因此,
\begin{align*}
\|T(u)-T(v)\|\leqslant c\cdot\|u-v\|.
\end{align*}
\end{proof}

在直线上的 Lebesgue 测度的研究中, 我们已观察到连续函数一般不将 Lebesgue 可测集映到 Lebesgue 可测集. 然而, 连续 Lipschitz 映射将 Lebesgue 可测集映到 Lebesgue 可测集.

\begin{proposition}\label{proposition:ch20-lipschitz可测}
令映射 $\Psi:\mathbf{R}^{n}\to\mathbf{R}^{n}$ 为 Lipschitz 的. 若 $E$ 是 $\mathbf{R}^{n}$ 的 Lebesgue 可测子集, 则 $\Psi(E)$ 也是. 特别地, 线性算子 $T:\mathbf{R}^{n}\to\mathbf{R}^{n}$ 将 Lebesgue 可测集映到 Lebesgue 可测集.
\end{proposition}

\begin{proof}
$\mathbf{R}^{n}$ 的子集是紧的当且仅当它是闭有界的且连续函数将紧集映到紧集. 由于 $\Psi$ 是 Lipschitz 的, 它是连续的. 因此 $\Psi$ 将有界 $F_{\sigma}$ 集映到 $F_{\sigma}$ 集.

令 $E$ 为 $\mathbf{R}^{n}$ 的 Lebesgue 可测子集. 由于 $\mathbf{R}^{n}$ 是有界可测集的可数族的并, 我们可以假设 $E$ 是有界的. 根据推论\cref{corollary:ch20-可测集刻画}, $E=A\cup D$, 其中 $A$ 是 $\mathbf{R}^{n}$ 的 $F_{\sigma}$ 子集, 而 $D$ 的 Lebesgue 外测度为零. 我们刚观察到 $\Psi(A)$ 是 $F_{\sigma}$ 集. 因此, 为证明 $\Psi(E)$ 是 Lebesgue 可测的, 仅须证明 $\Psi(D)$ 有零 Lebesgue 外测度.

令 $c>0$ 使得
\begin{align*}
\text{对所有 }u,v\in\mathbf{R}^{n},\quad\|\Psi(u)-\Psi(v)\|\leqslant c\|u-v\|.
\end{align*}
存在仅依赖于 $c$ 与 $n$ 的常数 $c'$(见习题 24)使得对 $\mathbf{R}^{n}$ 中的任何区间 $I$,
\begin{align}\label{align:ch20-lipschitz外测度}
\mu_{n}^{*}(\Psi(I))\leqslant c'\cdot\operatorname{vol}(I).
\end{align}
令 $\varepsilon>0$. 由于 $\mu_{n}^{*}(D)=0$, 存在 $\mathbf{R}^{n}$ 中覆盖 $D$ 的区间的可数族 $\{I^{k}\}_{k=1}^{\infty}$ 使得 $\sum_{k=1}^{\infty}\operatorname{vol}(I^{k})<\varepsilon$. 则 $\{\Psi(I^{k})\}_{k=1}^{\infty}$ 是 $\Psi(D)$ 的可数覆盖. 因此根据估计式\eqref{align:ch20-lipschitz外测度}与外测度的可数单调性,
\begin{align*}
\mu_{n}^{*}(\Psi(D))\leqslant\sum_{k=1}^{\infty}\mu_{n}^{*}(\Psi(I_{k}))\leqslant\sum_{k=1}^{\infty}c'\cdot\operatorname{vol}(I_{k})<c'\cdot\varepsilon.
\end{align*}
由于这对所有 $\varepsilon>0$ 成立, 它对 $\varepsilon=0$ 也成立.
\end{proof}

\begin{corollary}\label{corollary:ch20-复合可测}
令函数 $f:\mathbf{R}^{n}\to\mathbf{R}$ 关于 Lebesgue 测度可测, 而算子 $T:\mathbf{R}^{n}\to\mathbf{R}^{n}$ 是线性与可逆的. 则复合 $f\circ T:\mathbf{R}^{n}\to\mathbf{R}$ 关于 Lebesgue 测度也是可测的.
\end{corollary}

\begin{proof}
令 $\mathcal{O}$ 为 $\mathbf{R}$ 的开子集. 我们必须证明 $(f\circ T)^{-1}(\mathcal{O})$ 是 Lebesgue 可测的. 然而,
\begin{align*}
(f\circ T)^{-1}(\mathcal{O})=T^{-1}(f^{-1}(\mathcal{O})).
\end{align*}
函数 $f$ 是可测的, 因此集合 $f^{-1}(\mathcal{O})$ 是可测的. 另一方面, 映射 $T^{-1}$ 是线性的, 根据前一个命题, 它将 Lebesgue 可测集映到 Lebesgue 可测集. 因此 $(f\circ T)^{-1}(\mathcal{O})$ 是 Lebesgue 可测的.
\end{proof}

我们将从维数 $n=1$ 与 $n=2$ 开始证明关于 $\mathbf{R}^{n}$ 上的 Lebesgue 积分的一般换元公式.

\begin{proposition}\label{proposition:ch20-一维换元}
令 $f:\mathbf{R}\to\mathbf{R}$ 在 $\mathbf{R}$ 上关于一维 Lebesgue 测度可积. 若 $\alpha,\beta\in\mathbf{R}$, $\alpha\neq 0$, 则
\begin{align}\label{align:ch20-一维换元公式}
\int_{\mathbf{R}}f\,d\mu_{1}=|\alpha|\cdot\int_{\mathbf{R}}f(\alpha x)\,d\mu_{1}(x),\quad\int_{\mathbf{R}}f\,d\mu_{1}=\int_{\mathbf{R}}f(x+\beta)\,d\mu_{1}(x).
\end{align}
\end{proposition}

\begin{proof}
根据积分的线性, 可以假设 $f$ 是非负的. 用递增的简单可积函数序列逼近 $f$, 因此可用单调收敛定理将证明归结为 $f$ 是有限 Lebesgue 测度集的特征函数的情形. 对这样的函数, 公式是显然的.
\end{proof}

\begin{proposition}\label{proposition:ch20-二维换元}
令 $f:\mathbf{R}^{2}\to\mathbf{R}$ 在 $\mathbf{R}^{2}$ 上关于 Lebesgue 测度 $\mu_{2}$ 可积而 $c\neq 0$ 为实数. 对所有 $(x,y)\in\mathbf{R}^{2}$, 定义 $\varphi:\mathbf{R}^{2}\to\mathbf{R}$, $\psi:\mathbf{R}^{2}\to\mathbf{R}$ 和 $\eta:\mathbf{R}^{2}\to\mathbf{R}$ 为
\begin{align*}
\varphi(x,y)=f(y,x),\quad\psi(x,y)=f(x,x+y)\text{ 和 }\eta(x,y)=f(cx,y).
\end{align*}
则 $\varphi$, $\psi$ 与 $\eta$ 在 $\mathbf{R}^{2}$ 上关于 Lebesgue 测度 $\mu_{2}$ 可积. 此外,
\begin{align*}
\int_{\mathbf{R}^{2}}f\,d\mu_{2}=\int_{\mathbf{R}^{2}}\varphi\,d\mu_{2}=\int_{\mathbf{R}^{2}}\psi\,d\mu_{2}
\end{align*}
且
\begin{align*}
\int_{\mathbf{R}^{2}}f\,d\mu_{2}=|c|\cdot\int_{\mathbf{R}^{2}}\eta\,d\mu_{2}.
\end{align*}
\end{proposition}

\begin{proof}
我们从推论\cref{corollary:ch20-复合可测}推出函数 $\varphi$, $\psi$ 与 $\eta$ 都是 $\mu_{2}$ 可测的. 由于积分是线性的, 我们可以假设 $f$ 是非负的. 我们比较 $f$ 与 $\varphi$ 的积分且将其余两个积分的比较留作练习. 由于 $f$ 是 $\mu_{2}$ 可测的, 我们从定理\cref{theorem:ch20-富比尼Lebesgue}表述的 Fubini 定理推出,
\begin{align*}
\int_{\mathbf{R}^{2}}f\,d\mu_{2}=\int_{\mathbf{R}}\left[\int_{\mathbf{R}}f(x,y)\,d\mu_{1}(x)\right]d\mu_{1}(y).
\end{align*}
然而, 根据函数 $\varphi$ 的定义, 对几乎所有 $y\in\mathbf{R}$,
\begin{align*}
\int_{\mathbf{R}}f(x,y)\,d\mu_{1}(x)=\int_{\mathbf{R}}\varphi(y,x)\,d\mu_{1}(x).
\end{align*}
因此
\begin{align*}
\int_{\mathbf{R}}\left[\int_{\mathbf{R}}f(x,y)\,d\mu_{1}(x)\right]d\mu_{1}(y)=\int_{\mathbf{R}}\left[\int_{\mathbf{R}}\varphi(y,x)\,d\mu_{1}(x)\right]d\mu_{1}(y).
\end{align*}
由于 $\varphi$ 是非负的与 $\mu_{2}$ 可测的, 我们从定理\cref{theorem:ch20-富比尼Lebesgue}表述的 Tonelli 定理推出,
\begin{align*}
\int_{\mathbf{R}}\left[\int_{\mathbf{R}}\varphi(y,x)\,d\mu_{1}(x)\right]d\mu_{1}(y)=\int_{\mathbf{R}^{2}}\varphi\,d\mu_{2}.
\end{align*}
因此
\begin{align*}
\int_{\mathbf{R}^{2}}f\,d\mu_{2}=\int_{\mathbf{R}^{2}}\varphi\,d\mu_{2}.
\end{align*}
\end{proof}

目前为止我们所用到的线性映射仅有的分析性质是, 这样的映射是 Lipschitz 的. 我们现在需要线性代数的两个结果. 第一个是每个算子 $T\in GL(n,\mathbf{R})$ 可表示为以下三种初等类型的线性算子的复合:
\begin{description}
\item[类型 1] $T(e_{j})=ce_{j}$ 且对 $k\neq j$, $T(e_{k})=e_{k}$.
\item[类型 2] $T(e_{j})=e_{j+1}$, $T(e_{j+1})=e_{j}$ 且对 $k\neq j$, $j+1$, $T(e_{k})=e_{k}$.
\item[类型 3] $T(e_{j})=e_{j}+e_{j+1}$ 且对 $k\neq j$, $T(e_{k})=e_{k}$.
\end{description}
每个可逆线性算子可表示为初等算子的复合是根据线性算子的矩阵的性质做出的断言: 每个可逆 $n\times n$ 矩阵可从单位矩阵的行运算得到.

我们需要的线性算子的第二个性质是: 对每个线性算子 $T:\mathbf{R}^{n}\to\mathbf{R}^{n}$, 存在与之相应的称为 $T$ 的行列式(记为 $\det T$)的实数, 它具有以下三个性质:
\begin{enumerate}[(i)]
\item 对任何两个线性算子 $T$, $S:\mathbf{R}^{n}\to\mathbf{R}^{n}$
\begin{align}\label{align:ch20-行列式乘性}
\det(S\circ T)=\det S\cdot\det T.
\end{align}
\item 若 $T$ 是类型 1, $\det T=c$; 若 $T$ 是类型 2, $\det T=-1$; 若 $T$ 是类型 3, $\det T=1$.
\item 若 $T(e_{n})=e_{n}$ 且 $T$ 将子空间 $\{x\in\mathbf{R}^{n}\mid x=(x_{1},x_{2},\cdots,x_{n-1},0)\}$ 映到其自身, 则 $\det T=\det T'$, 其中 $T':\mathbf{R}^{n-1}\to\mathbf{R}^{n-1}$ 是 $T$ 在 $\mathbf{R}^{n-1}$ 上的限制.
\end{enumerate}

\begin{theorem}\label{theorem:ch20-换元公式}
令线性算子 $T:\mathbf{R}^{n}\to\mathbf{R}^{n}$ 可逆. 假定函数 $f:\mathbf{R}^{n}\to\mathbf{R}$ 在 $\mathbf{R}^{n}$ 上关于 Lebesgue 测度可积. 则复合 $f\circ T:\mathbf{R}^{n}\to\mathbf{R}$ 也在 $\mathbf{R}^{n}$ 上关于 Lebesgue 测度可积且
\begin{align}\label{align:ch20-一般换元}
\int_{\mathbf{R}^{n}}f\circ T\,d\mu_{n}=\frac{1}{|\det T|}\cdot\int_{\mathbf{R}^{n}}f\,d\mu_{n}.
\end{align}
\end{theorem}

\begin{proof}
积分是线性的. 我们因此假定 $f$ 是非负的. 根据行列式的可乘性与将可逆算子分解为初等算子的复合, 我们可以假定 $T$ 是初等的. $n=1$ 的情形被\eqref{align:ch20-一维换元公式}涵盖. 根据命题\cref{proposition:ch20-二维换元}, \eqref{align:ch20-一般换元}对 $n=2$ 也成立. 我们现在用归纳法. 假设我们已对 $m\geqslant 2$ 证明了\eqref{align:ch20-一般换元}且考虑 $n=m+1$ 的情形. 由于 $T$ 是初等的且 $n\geqslant 3$, 或者(i) $T(e_{n})=e_{n}$ 且 $T$ 将子空间 $\{x\in\mathbf{R}^{n}\mid x=(x_{1},x_{2},\cdots,x_{n-1},0)\}$ 映到其自身, 或者(ii) $T(e_{1})=e_{1}$ 且 $T$ 将子空间 $\{x\in\mathbf{R}^{n}\mid x=(0,x_{2},\cdots,x_{n})\}$ 映到其自身. 我们考虑情形(i)且将类似的情形(ii)留作练习. 令 $T'$ 为由 $T$ 在 $\mathbf{R}^{n-1}$ 上诱导的算子. 观察到 $|\det T'|=|\det T|$. 我们现在再次用命题\cref{proposition:ch20-二维换元}的证明中的方法. 函数 $f\circ T$ 是 $\mu_{n}$ 可测的. 因此我们从定理\cref{theorem:ch20-富比尼Lebesgue}中对 Lebesgue 测度阐述的 Fubini 定理与 Tonelli 定理以及\eqref{align:ch20-一般换元}对 $m=n-1$ 成立推出,
\begin{align*}
\int_{\mathbf{R}^{n}}f\circ T\,d\mu_{n}&=\int_{\mathbf{R}}\left[\int_{\mathbf{R}^{n-1}}f\circ T(x_{1},x_{2},\cdots,x_{n})\,d\mu_{n-1}(x_{1},\cdots,x_{n-1})\right]d\mu_{1}(x_{n})\\
&=\int_{\mathbf{R}}\left[\int_{\mathbf{R}^{n-1}}f(T'(x_{1},\cdots,x_{n-1}),x_{n})\,d\mu_{n-1}(x_{1},\cdots,x_{n-1})\right]d\mu_{1}(x_{n})\\
&=\frac{1}{|\det T'|}\int_{\mathbf{R}}\left[\int_{\mathbf{R}^{n-1}}f(x_{1},x_{2},\cdots,x_{n})\,d\mu_{n-1}(x_{1},\cdots,x_{n-1})\right]d\mu_{1}(x_{n})\\
&=\frac{1}{|\det T|}\int_{\mathbf{R}^{n}}f\,d\mu_{n}.
\end{align*}
\end{proof}

\begin{corollary}\label{corollary:ch20-线性变换测度}
令线性算子 $T:\mathbf{R}^{n}\to\mathbf{R}^{n}$ 是可逆的. 则对 $\mathbf{R}^{n}$ 的每个 Lebesgue 可测子集 $E$, $T(E)$ 是 Lebesgue 可测的且
\begin{align}\label{align:ch20-变换测度}
\mu_{n}(T(E))=|\det(T)|\cdot\mu_{n}(E).
\end{align}
\end{corollary}

\begin{proof}
我们假设 $E$ 是有界的而将无界的情形留作练习. 由于 $T$ 是 Lipschitz 的, $T(E)$ 是有界的. 我们从命题\cref{proposition:ch20-lipschitz可测}推出 $T(E)$ 是 Lebesgue 可测的, 且由于它是有界的, 它具有有限 Lebesgue 测度. 因此函数 $f=\chi_{T(E)}$ 在 $\mathbf{R}^{n}$ 上关于 Lebesgue 测度是可积的. 观察到 $f\circ T=\chi_{E}$. 因此,
\begin{align*}
\int_{\mathbf{R}^{n}}f\circ T\,d\mu_{n}=\mu_{n}(E)\text{ 且 }\int_{\mathbf{R}^{n}}f\,d\mu_{n}=\mu_{n}(T(E)).
\end{align*}
因此对 $f$ 的这个特定选取, \eqref{align:ch20-变换测度}从\eqref{align:ch20-一般换元}得到.
\end{proof}

谈到 $\mathbf{R}^{n}$ 的刚体运动, 我们意味着将 $\mathbf{R}^{n}$ 映上 $\mathbf{R}^{n}$ 且保持点之间的欧氏距离的映射 $\psi$, 即
\begin{align*}
\text{对所有 }u,v\in\mathbf{R}^{n},\quad\|\Psi(u)-\Psi(v)\|=\|u-v\|.
\end{align*}
Mazur 与 Ulam 的一个定理告诉我们, 每个刚体运动是线性刚体运动的常数扰动, 即存在 $\mathbf{R}^{n}$ 中的点 $x_{0}$ 与线性算子 $T:\mathbf{R}^{n}\to\mathbf{R}^{n}$ 使得对所有 $x\in\mathbf{R}^{n}$, $\Psi(x)=T(x)+x_{0}$, 其中 $T$ 是线性刚体运动. 然而, 由于线性刚体运动将原点映到原点, $T$ 保持范数, 即对所有 $u\in\mathbf{R}^{n}$,
\begin{align*}
\|Tu\|=\|u\|.
\end{align*}
因此以下极化等式(见习题 28)
\begin{align}\label{align:ch20-极化等式}
\text{对所有 }u,v\in\mathbf{R}^{n},\quad\langle u,v\rangle=\frac{1}{4}\{\|u+v\|^{2}-\|u-v\|^{2}\}
\end{align}
告诉我们线性刚体运动 $T$ 保持内积, 即
\begin{align*}
\text{对所有 }u,v\in\mathbf{R}^{n},\quad\langle T(u),T(v)\rangle=\langle u,v\rangle.
\end{align*}
该等式意味着 $T^{*}T=Id$. 从行列式的乘性以及事实 $\det T=\det T^{*}$, 我们得出 $|\det T|=1$. 因此根据 Lebesgue 测度的平移不变性(见习题 20)和\eqref{align:ch20-变换测度}, 我们有以下有趣的几何结果: 若欧氏空间上的映射保持点的距离不变, 则它保持 Lebesgue 测度.

\begin{corollary}\label{corollary:ch20-刚体运动}
令 $\Psi:\mathbf{R}^{n}\to\mathbf{R}^{n}$ 为刚体运动. 则对 $\mathbf{R}^{n}$ 的每个 Lebesgue 可测子集 $E$,
\begin{align*}
\mu_{n}(\Psi(E))=\mu_{n}(E).
\end{align*}
\end{corollary}

从 Lebesgue 外测度 $\mu_{n}^{*}$ 的定义得出, $\mathbf{R}^{n}$ 的子空间 $V=\{x\in\mathbf{R}^{n}\mid x=(x_{1},x_{2},\cdots,x_{n-1},0)\}$ 的 $n$ 维 Lebesgue 测度为零. 我们从\eqref{align:ch20-变换测度}推出 $\mathbf{R}^{n}$ 的任何真子空间 $W$ 的 $n$ 维 Lebesgue 测度为零, 由于它可被 $GL(n,\mathbf{R})$ 中的算子映到 $V$ 的子空间. 由此得出若线性算子 $T:\mathbf{R}^{n}\to\mathbf{R}^{n}$ 不是可逆的, 则由于它的值域落在维数小于 $n$ 的子空间, 它将 $\mathbf{R}^{n}$ 的每个子集映到 $n$ 维 Lebesgue 测度为零的集合. 这可重述为\eqref{align:ch20-变换测度}对不可逆算子 $T$ 仍然成立.

\subsection*{习题}

\begin{exercise}\label{exercise:ch20-15}
考虑三角形 $\Delta=\{(x,y)\in\mathbf{R}^{2}\mid 0\leqslant x\leqslant a,\ 0\leqslant y\leqslant[b/a]x\}$. 用矩形的有限族覆盖 $\Delta$ 且用测度的连续性, 确定 $\Delta$ 的 Lebesgue 测度.
\end{exercise}

\begin{exercise}\label{exercise:ch20-16}
令 $[a,b]$ 为实数的闭有界区间. 假定 $f:[a,b]\to\mathbf{R}$ 是有界与 Lebesgue 可测的. 证明 $f$ 的图像关于平面的 Lebesgue 测度为零. 将此推广到多个实变量的有界实值函数.
\end{exercise}

\begin{exercise}\label{exercise:ch20-17}
实数的每个开子集是可数个不交开区间的并. 平面的开子集 $\{(x,y)\in\mathbf{R}^{2}\mid 0<x,\ y<1\}$ 是可数个不交开球的并吗?
\end{exercise}

\begin{exercise}\label{exercise:ch20-18}
证明不等式\eqref{align:ch20-一维整点估计}.
\end{exercise}

\begin{exercise}\label{exercise:ch20-19}
证明集合等式\eqref{align:ch20-开集方体分解}.
\end{exercise}

\begin{exercise}\label{exercise:ch20-20}
令 $E\subseteq\mathbf{R}^{n}$ 而 $z\in\mathbf{R}^{n}$.
\begin{enumerate}[(i)]
\item 证明: 若 $E$ 是开的, $E+z$ 也是开的.
\item 证明: 若 $E$ 是 $G_{\delta}$ 集, $E+z$ 也是 $G_{\delta}$ 集.
\item 证明 $\mu_{n}^{*}(E+z)=\mu_{n}^{*}(E)$.
\item 证明 $E$ 是 $\mu_{n}$ 可测的当且仅当 $E+z$ 是 $\mu_{n}$ 可测的.
\end{enumerate}
\end{exercise}

\begin{exercise}\label{exercise:ch20-21}
对每个自然数 $n$, 证明 $\mathbf{R}^{n}$ 的每个具有正外 Lebesgue 测度的子集包含一个不是 Lebesgue 可测的子集.
\end{exercise}

\begin{exercise}\label{exercise:ch20-22}
对每个自然数 $n$, 证明存在 $\mathbf{R}^{n}$ 的子集, 它不是 Borel 集但却是 $\mu_{n}$ 可测的.
\end{exercise}

\begin{exercise}\label{exercise:ch20-23}
若\eqref{align:ch20-变换测度}对 $\mathbf{R}^{n}$ 中的每个区间成立, 用 Carath\'eodory-Hahn 定理的唯一性断言直接证明它对 $\mathbf{R}^{n}$ 的每个可测子集也成立.
\end{exercise}

\begin{exercise}\label{exercise:ch20-24}
令 $\Psi:\mathbf{R}^{n}\to\mathbf{R}^{n}$ 为具有 Lipschitz 常数 $c$ 的 Lipschitz 映射. 证明存在常数 $c'$ 使得估计式\eqref{align:ch20-lipschitz外测度}成立.
\end{exercise}

\begin{exercise}\label{exercise:ch20-25}
证明两个半环的笛卡儿积是半环. 以此为基础用归纳法证明 $\mathbf{R}^{n}$ 的区间族是半环.
\end{exercise}

\begin{exercise}\label{exercise:ch20-26}
若函数 $f:[0,1]\times[0,1]\to\mathbf{R}$ 关于每个变量分别连续, 则它关于 Lebesgue 测度 $\mu_{2}$ 可测.
\end{exercise}

\begin{exercise}\label{exercise:ch20-27}
令 $g:\mathbf{R}\to\mathbf{R}$ 为 $\mathbf{R}$ 映上 $\mathbf{R}$ 的映射使得存在常数 $c>0$,
\begin{align*}
\text{对所有 }u,v\in\mathbf{R},\ |g(u)-g(v)|\geqslant c\cdot|u-v|.
\end{align*}
证明: 若 $f:\mathbf{R}\to\mathbf{R}$ 是 Lebesgue 可测的, 则复合 $f\circ g:\mathbf{R}\to\mathbf{R}$ 也如此.
\end{exercise}

\begin{exercise}\label{exercise:ch20-28}
用内积的双线性形式证明\eqref{align:ch20-极化等式}.
\end{exercise}

\begin{exercise}\label{exercise:ch20-29}
令映射 $T:\mathbf{R}^{n}\to\mathbf{R}^{n}$ 为线性的. 定义 $c=\sup\{\|T(x)\|\mid\|x\|\leqslant 1\}$. 证明 $c$ 是 $T$ 的最小 Lipschitz 常数.
\end{exercise}

\begin{exercise}\label{exercise:ch20-30}
证明 $\mathbf{R}^{n}$ 的维数小于 $n$ 的子空间 $W$ 的 $n$ 维 Lebesgue 测度为零, 首先对子空间 $\{x\in\mathbf{R}^{n}\mid x_{n}=0\}$ 证明这个结论.
\end{exercise}

\begin{exercise}\label{exercise:ch20-31}
首先对有限测度集的特征函数, 接着对有限测度集外消失的简单函数, 最后对一元实变量的非负可积函数证明换元公式\eqref{align:ch20-一维换元公式}.
\end{exercise}

\begin{exercise}\label{exercise:ch20-32}
对 $\mathbf{R}$ 的子集 $E$, 定义
\begin{align*}
\sigma(E)=\{(x,y)\in\mathbf{R}^{2}\mid x-y\in E\}.
\end{align*}
\begin{enumerate}[(i)]
\item 若 $E$ 是 $\mathbf{R}$ 的 Lebesgue 可测子集, 证明 $\sigma(E)$ 是 $\mathbf{R}^{2}$ 的 Lebesgue 可测子集. (提示: 考虑 $E$ 是开, $E$ 是 $G_{\delta}$, $E$ 的测度为零, $E$ 是可测的情形.)
\item 若 $f$ 是 $\mathbf{R}$ 上的 Lebesgue 可测函数, 证明函数 $F(x,y)=f(x-y)$ 是 $\mathbf{R}^{2}$ 上的 Lebesgue 可测函数.
\item 若 $f$ 和 $g$ 属于 $L^{1}(\mathbf{R},\mu_{1})$, 证明对 $\mathbf{R}$ 中的几乎所有 $x$, 由 $\varphi(y)=f(x-y)g(y)$ 给出的函数 $\varphi$ 属于 $L^{1}(\mathbf{R},\mu_{1})$. 若我们记它的积分为 $h(x)$, 证明 $h$ 是可积的且
\begin{align*}
\int_{\mathbf{R}}|h|\,d\mu_{1}\leqslant\int_{\mathbf{R}}|f|\,d\mu_{1}\cdot\int_{\mathbf{R}}|g|\,d\mu_{1}.
\end{align*}
\end{enumerate}
\end{exercise}

\begin{exercise}\label{exercise:ch20-33}
令 $f$ 和 $g$ 为 $L^{1}(\mathbf{R},\mu_{1})$ 中的函数, 定义 $\mathbf{R}$ 上的 $f*g$ 为
\begin{align*}
(f*g)(y)=\int_{\mathbf{R}}f(y-x)g(x)\,d\mu_{1}(x).
\end{align*}
\begin{enumerate}[(i)]
\item 证明 $f*g=g*f$.
\item 证明对每个 $h\in L^{1}(\mathbf{R},\mu_{1})$, $(f*g)*h=f*(g*h)$.
\end{enumerate}
\end{exercise}

\begin{exercise}\label{exercise:ch20-34}
令 $f$ 为在 $\mathbf{R}$ 上关于 $\mu_{1}$ 可积的非负函数. 证明
\begin{align*}
\mu_{2}\{(x,y)\in\mathbf{R}^{2}\mid 0\leqslant y\leqslant f(x)\}=\mu_{2}\{(x,y)\in\mathbf{R}^{2}\mid 0<y<f(x)\}=\int_{\mathbf{R}}f(x)\,dx.
\end{align*}
对每个 $t\geqslant 0$, 定义 $\varphi(t)=\mu_{1}\{x\in\mathbf{R}\mid f(x)\geqslant t\}$. 证明 $\varphi$ 是递减函数且
\begin{align*}
\int_{0}^{\infty}\varphi(t)\,d\mu_{n}(t)=\int_{\mathbf{R}}f(x)\,d\mu_{1}(x).
\end{align*}
\end{exercise}

\end{document}